\documentclass[10pt,a4paper]{amsart}
\usepackage[margin=19mm]{geometry}
\usepackage{amsmath,amssymb,mathtools}
\usepackage[scr=rsfs,cal=euler]{mathalfa}
\usepackage{xfrac}
\usepackage[unicode,hidelinks,bookmarks=false]{hyperref}
\newcommand{\ie}{i.e.\ }
\newcommand{\cf}{cf.\ }
\newcommand{\resp}{respectively\ }
\newcommand{\Subsection}{\subsection}
\providecommand{\nobreakdash}{\nobreak\hbox{-}}
\setlength{\emergencystretch}{2em}
\multlinegap0pt
\usepackage{amssymb}
\usepackage{algorithm}\usepackage{algpseudocode}
\newtheorem{theorem}{Theorem}
\def\C{\mathbb{C}}
\def\R{\mathbb{R}}
\def\Z{\mathbb{Z}}
\newcommand{\be}{\mathbf e}
\newcommand{\bmu}{{\boldsymbol{\mu}}}
\newcommand{\bu}{\mathbf u}
\newcommand{\bw}{\mathbf w}
\def\ii{\mathrm{i}}
\newcommand{\uno}{\mathbf 1}
\newcommand{\zero}{\mathbf 0}
\title{Perfect $(s,r)$-state transfer}
\author{Stephen Kirkland, Hermie Monterde,, Sarah Plosker}
\date{Source: arXiv:2609.10713v1 (CC BY 4.0)}
\begin{document}
\maketitle

\noindent\emph{Source and licence.} Perfect $(s,r)$-state transfer. Version arXiv:2609.10713v1 (CC BY 4.0), distributed by arXiv under the Creative Commons Attribution 4.0 International licence, http://creativecommons.org/licenses/by/4.0/, as recorded in the arXiv licence metadata for this exact version and shown on the abstract page https://arxiv.org/abs/2609.10713v1. Full licence notice: \texttt{licenses/arxiv-2609.10713-CC-BY-4.0.txt}.\par\smallskip

\noindent\textit{Editorial scope.} Counted unit: the article's theorem Thm:kmnPSTc1 (source0.tex lines 309--348). The article's complete original proof of it is delivered with it (source0.tex lines 350--442, one \texttt{proof} environment of 93 lines). No mathematics is written, repaired or re-derived: the statement and the proof are the article's own lines, in the article's own notation, and no retained line is substituted or re-flowed.  Retained uncounted as the source-local context the proof needs: lines 298--307.  Nothing was omitted from any retained span, so the package carries no omission record.  No retained line contains a citation, so the wrapper carries no bibliography and no bibliography entry.  No retained line contains a figure environment, an image-inclusion command or drawing code, so no image is delivered and the package declares no asset dependency.  Editorial formatting only, in addition: an AMS \texttt{amsart} preamble that restates the article's own declarations and macros used by the retained lines, namely the environments \texttt{theorem} on separate counters, together with the standard packages those lines need.  The retained environments print in this document as Theorem 1 in order of appearance, whereas the article prints the same items with the numbering of its own full document.  The complete native source of the article as distributed in the arXiv e-print for this exact version is bundled unchanged as \texttt{sources/209942/source0.tex}.
\begin{equation}
\label{Eq:completebip}
U_A(t)=\begin{bmatrix}
I_a-\frac{1}{a}J&0 \\ 0&I_b-\frac{1}{b}J
\end{bmatrix}+\frac{e^{\ii t\sqrt{ab}}}{2ab}\begin{bmatrix}
bJ&\sqrt{ab}J \\ \sqrt{ab}J&aJ
\end{bmatrix}+\frac{e^{-it\sqrt{ab}}}{2ab}\begin{bmatrix}
bJ&-\sqrt{ab}J \\ -\sqrt{ab}J&aJ
\end{bmatrix}.
\end{equation}

\begin{theorem}
\label{Thm:kmnPSTc1}
Let $V(K_{a,b})=B_1\cup B_2$, where $|B_1|=a$. Suppose $\bu=\be_u+s\be_v$ for some $s\in\C\backslash\{0\}$ and $\bmu\in\C^n$ is a vector with two nonzero entries. PST occurs from $\bu$ to $\bmu$ in $K_{a,b}$ relative to $A$ at time $\tau$ if and only if one of the following conditions hold.
\begin{enumerate}
\item $u,v\in B_1$, and one of the conditions below holds.
\begin{enumerate}

\item $*$ $a=2$, $b\geq 1$, $s\neq 1$, $\bmu=-(\be_v+s\be_u)$ and $\tau=\frac{\pi}{\sqrt{2b}}$. 

\item $*$ $a=4$, $b\geq 1$, $s=1$, $\bmu=\be_w+\be_x$ where $B_1=\{u,w,v,x\}$, and $\tau=\frac{\pi}{2\sqrt{b}}$. 

\item $*$ $a=3$, $b\geq 1$, $s=\frac{1}{2}$, $\bmu=\be_w+\frac{1}{2}\be_x$ where $B_1=\{u,w,v\}$, and $\tau=\frac{\pi}{\sqrt{3b}}$. 

\item $\dagger$ $a=2$, $b=1$, $s=\cot^2(\tau/\sqrt{2})
$, $\bmu=(s-1)\be_v+\frac{1+s}{\sqrt{2}}\ii\sin(\tau\sqrt{2})\be_w$ where $B_2=\{w\}$, and $\tau\in\R$.
\item $\dagger$ $a=2$, $b=1$, $s=\tan^2(\tau/\sqrt{2})$, $\bmu=(1-s)\be_u+\frac{1+s}{\sqrt{2}}\ii\sin(\tau\sqrt{2})\be_w$ where $B_2=\{w\}$, and $\tau\in\R$.

\item $*$ $a=3$, $b=1$, $s=1$, $(\bmu,\tau)\in\{\big(-\be_x+\ii\be_w,\frac{2\pi(3k+1)}{3\sqrt{3}}\big),\big(-\be_x-\ii\be_w,\frac{2\pi(3k+2)}{3\sqrt{3}}\big)\}$ where $B_1=\{u,v,x\}$, $B_2=\{w\}$ and $k\in\Z$.

\item $*$ $a=2$, $b=2$, $s=1$, and $\bmu=\ii(\be_w+\be_x)$ where $B_2=\{w,x\}$, and $\tau=\frac{\pi}{4}$.
\end{enumerate}

\item $u\in B_1$, $v\in B_2$ and one of the conditions below hold.
\begin{enumerate}
\item $*$ $a=1$, $b=2$, $s\in\C\backslash\{0\}$, $\bmu=\be_u+s\be_w$ where $B_2=\{v,w\}$, and $\tau=\frac{\pi}{\sqrt{2}}$.
\item $*$ $a=b=2$, $s\in\C\backslash\{0\}$, $\bmu=\be_w+s\be_x$ where $B_1=\{u,w\}$ and $B_2=\{v,x\}$, and $\tau=\frac{\pi}{2}$.
\item $*$ $a=1$, $b\geq 1$, $s=\ii\sqrt{b}\cot(\tau\sqrt{b}/2)
$, $\bmu=-\be_u+s\be_v$, and $\tau\in\R$. 
\item $\dagger$ $a=1$, $b=2$, $s=-\bigg(\frac{1}{\sqrt{2}\csc(\tau\sqrt{2})-\sqrt{1/2}\tan(\tau/\sqrt{2})}\bigg)\ii $, $\bmu=\left(\cos( \tau\sqrt{2})+\ii \frac{s}{\sqrt{2}}\sin(\tau\sqrt{2})\right)\be_u -s\be_w$
where $B_2=\{v,w\}$, and $\tau\in\R$.
\item $\dagger$ $a=1$, $b=2$, $s=-\sqrt{2}\cot(\tau\sqrt{2})$, %$s=\sqrt{2}\big(\csc(\tau\sqrt{2})-\tan(\tau/\sqrt{2})\big)\ii$,  
$B_2=\{v,w\}$, $\bmu=\left(\ii\frac{1}{\sqrt{2}}\sin(\tau\sqrt{2})+\frac{s}{2}(\cos(\tau\sqrt{2})+1)\right)\be_v+\left((\ii\sin(\tau\sqrt{2})\frac{1}{\sqrt{2}}+\frac{s}{2}(\cos(\tau\sqrt{2})-1)\right)\be_w$, and $\tau\in\R$.
\item $*$ $a=1$, $b=3$, $(s,\bmu,\tau)\in\left\{\big(1,\ii(\be_w+\be_x),\frac{2\pi(3k+1)}{3\sqrt{3}}\big),\big(1,-\ii(\be_w+\be_x),\frac{2\pi(3k+2)}{3\sqrt{3}}\big)\right\}$, where $B_2=\{v,w,x\}$ and $k\in \Z$.


\item $*$ $a=2$, $b\geq 1$, $s=\ii \sqrt{b/2} \cot(\tau\sqrt{b/2})$, $\bmu=-\be_w+s\be_v$ where $B_1=\{u,w\}$, and $\tau\in\R$.
\end{enumerate}
\end{enumerate}
Moreover, the above statements all hold if we interchange the roles of $a$ and $b$, as well as $B_1$ and $B_2$.
\end{theorem}

\begin{proof}
First, let $u,v\in B_1$ so that $a\geq 2$. Then $s\neq -1$, otherwise $\bu$ is a fixed state. Letting $\bu_1=\be_u+s\be_v\in\C^a$ and $\bu_2=\zero_b$, (\ref{Eq:completebip}) gives us
\begin{equation*}
\bmu=U_A(\tau)\begin{bmatrix}
\bu_1\\ \bu_2\end{bmatrix}=\begin{bmatrix}
\be_u+s\be_v+\frac{1+s}{a}\big(\cos( \tau\sqrt{ab})-1\big)\uno_a \\ \frac{1+s}{\sqrt{ab}}\ii\sin( \tau\sqrt{ab})\uno_b
\end{bmatrix}. 
\end{equation*}
We have two cases. 

\noindent \textbf{Case 1.} Suppose $\sin( \tau\sqrt{ab})=0$. That is, $\tau=\frac{\ell\pi}{\sqrt{ab}}$ for any integer $\ell$. If $\ell$ is even, then $\bu=\bmu$ is periodic at $\tau$. If $\ell$ is odd, 
$\bmu$ has two nonzero entries if and only if $a=2$ and $s\neq -1$, or $(s,a)\in\{(1,4),(\frac{1}{2},3)\}$. If $a=2$ and $s\neq -1$, then $\bmu=-(\be_v+s\be_u)$, and so 1a holds. If $(s,a)=(1,4)$, then PST occurs between $\be_u+\be_v$ and $\be_w+\be_x$ at $\frac{\pi}{2\sqrt{b}}$, where $B_1=\{u,w,v,x\}$. So 1b holds. If $(s,a)=(\frac{1}{2},3)$, then PST occurs between $\be_u+\frac{1}{2}\be_v$ and $\be_w+\frac{1}{2}\be_v$ at $\frac{\pi}{\sqrt{3b}}$, where $B_1=\{u,w,v\}$. This proves 1c. 

\noindent \textbf{Case 2.} Suppose that $\sin( \tau\sqrt{ab})\neq 0$. Since $s\neq -1$ and $\bmu$ has exactly two nonzero entries, we have $b\leq 2$. If $b=1$, then
\begin{equation*}
\bmu=\begin{bmatrix}
\be_u+s\be_v+\frac{1+s}{a}(\cos(\tau\sqrt{a})-1)\uno_a \\ \frac{1+s}{\sqrt{a}}\ii\sin(\tau\sqrt{a})
\end{bmatrix}
\end{equation*}
has two nonzero entries if and only if (i) $a=2$ and $\frac{1+s}{2}(\cos(\tau\sqrt{2})-1)\in\{-1,-s\}$ or (ii) $a=3$, $s=1$ and $\tau=\frac{2\pi(3k+\ell)}{3\sqrt{3}}$, where $\ell\in\{1,2\}$ and $k\in\Z$. Note that (i) yields the values 
\begin{equation*}
s=\dfrac{1+\cos(\tau\sqrt{2})}{1-\cos(\tau\sqrt{2})}=\cot^2(\tau/\sqrt{2})\quad  \text{and} \quad s=\dfrac{1-\cos(\tau\sqrt{2})}{1+\cos(\tau\sqrt{2})}=\tan^2(\tau/\sqrt{2})
%=\sec(\tau/\sqrt{2})-1,
\end{equation*}
and the vectors $\bmu=(s-1)\be_v+\frac{1+s}{\sqrt{2}}\ii\sin(\tau\sqrt{2})\be_w$ and $\bmu=(1-s)\be_u+\frac{1+s}{\sqrt{2}}\ii\sin(\tau\sqrt{2})\be_w$,
respectively, where $B_2=\{w\}$. This proves 1d and 1e. Meanwhile, (ii) gives us $\bmu=-\be_x+\ii\be_w$ whenever $\ell=1$ and $\bmu=-\be_x-\ii\be_w$ whenever $\ell=2$, where $B_1=\{u,v,x\}$ and $B_2=\{w\}$. This establishes 1f. On the other hand, if $b=2$, then $\bmu$ has exactly two nonzero entries if and only if $\be_u+s\be_v+\frac{1+s}{a}(\cos(\tau\sqrt{2a})-1)\uno=0$. Equivalently, $a=2$, $s=1$, $\tau=\frac{\pi}{4}$ and $\bmu=\ii (\be_w+\be_x)$, where $B_2=\{w,x\}$. Thus, 1g holds.

Next, suppose that $u\in B_1$ and $v\in B_2$. We again use (\ref{Eq:completebip}) to obtain an expression for $\bmu=U_A(t)\begin{bmatrix}
\bu_1\\ \bu_2\end{bmatrix}$, where $\bu_1=\be_u$ and $\bu_2=s\be_v$. One checks that
\begin{equation}
\label{Eq:pf}
\bmu=\begin{bmatrix}
\be_u+\left(\big(\cos( \tau\sqrt{ab})-1\big)\frac{1}{a}+\ii s\sin(\tau\sqrt{ab})\frac{1}{\sqrt{ab}}\right)\uno_a \\ s\be_v+\left(\ii\sin(\tau\sqrt{ab})\frac{1}{\sqrt{ab}}+s(\cos(\tau\sqrt{ab})-1)\frac{1}{b}\right)\uno_b
\end{bmatrix}. 
\end{equation}
We again have two cases.

\noindent \textbf{Case 1.} Assume $\sin( \tau\sqrt{ab})=0$. That is, $\tau=\frac{\ell\pi}{\sqrt{ab}}$ for any integer $\ell$. If $\ell$ is even, then $\bu=\bmu$ is again periodic at $\tau$. If $\ell$ is odd, then 
$\bmu=\begin{bmatrix}
\be_u-\frac{2}{a}\uno_a \\ s\big(\be_v-\frac{2}{b}\uno_b\big).
\end{bmatrix}$ has two nonzero entries if and only if $(a,b)\in\{(1,2),(2,2)\}$. From this, 2a and 2b are immediate. 

\noindent \textbf{Case 2.} Next, suppose $\sin( \tau\sqrt{ab})\neq 0$. 
%If $\sin(\tau\sqrt{ab})=1$, then $\cos(\tau\sqrt{ab})=0$, and so
%\begin{center}
%$\bmu=\begin{bmatrix}
%\be_u+\left(-\frac{1}{a}+\ii s\frac{1}{\sqrt{ab}}\right)\uno_a \\ s\be_v+\left(\ii\frac{1}{\sqrt{ab}}-s\frac{1}{b}\right)\uno_b
%\end{bmatrix}. $
%\end{center}
Suppose we write (\ref{Eq:pf}) as $\bmu:=\begin{bmatrix}
\bw_1 \\ \bw_2\end{bmatrix}$ where $\bw_1\in\C^a$. We make the following observations for the vector $\bw_1$. First, we have $\bw_1=\be_u$ if and only if 
\begin{equation}
\label{eqkmn1}
s=-\frac{\frac{1}{a}(\cos(\tau\sqrt{ab})-1)}{\ii\sin(\tau\sqrt{ab})/\sqrt{ab}}=-\ii \tan(\tau\sqrt{ab}/2)\sqrt{b/a}.
\end{equation}
If $a\geq 2$, then $\bw_1$ has $a-1$ nonzero entries if and only if 
\begin{equation}
\label{eqkmn2}
s=-\frac{1+\frac{1}{a}(\cos(\tau\sqrt{ab})-1)}{\ii\sin(\tau\sqrt{ab})/\sqrt{ab}}=\bigg(\sqrt{ab}\csc(\tau\sqrt{ab})-\tan(\tau\sqrt{ab}/2)\sqrt{b/a}\bigg)\ii,
\end{equation}
in which case $\bw_1$ has all entries equal to $-1$
except for the $u$th entry which is equal to 0. Moreover, if $a=1$, then $\bw_1$ is a scalar equal to 0 if and only if $s$ satisfies (\ref{eqkmn2}). Next, we make the following observations for the vector $\bw_2$. Note that $\bw_2=s\be_v$ if and only if 
\begin{equation}
\label{eqkmn3}
s=-\frac{\ii\sin(\tau\sqrt{ab})/\sqrt{ab}}{\frac{1}{b}(\cos(\tau\sqrt{ab})-1)}=\ii \cot(\tau\sqrt{ab}/2)\sqrt{b/a}.
\end{equation}
\item If $b\geq 2$, then $\bw_2$ has $b-1$ nonzero entries if and only if 
\begin{equation}
\label{eqkmn4}
s=-\frac{\ii\sin(\tau\sqrt{ab})/\sqrt{ab}}{1+\frac{1}{b}(\cos(\tau\sqrt{ab})-1)}=-\bigg(\frac{1}{\sqrt{ab}\csc(\tau\sqrt{ab})-\tan(\tau\sqrt{ab}/2)\sqrt{a/b}}\bigg)\ii.
\end{equation}
in which case $\bw_2$ has all entries equal to $-s$ 
except for the $v$th entry which is equal to 0. Moreover, if $b=1$, then $\bw_2$ is a scalar equal to 0 if and only if $s$ satisfies (\ref{eqkmn4}).
Additionally, if $s$ is not equal to the values in (\ref{eqkmn1}) and (\ref{eqkmn2}), then all entries of $\bw_1$ are nonzero. Similarly, if $s$ is not equal to the values in (\ref{eqkmn3}) and (\ref{eqkmn4}), then all entries of $\bw_2$ are nonzero. Furthermore, (\ref{eqkmn1}) and (\ref{eqkmn2}) cannot simultaneously hold, (\ref{eqkmn3}) and (\ref{eqkmn4}) cannot simultaneously hold, and (\ref{eqkmn1}) and (\ref{eqkmn3}) cannot simultaneously hold (otherwise $\bmu$ is periodic). One also checks that (\ref{eqkmn2}) and (\ref{eqkmn3}) simultaneously hold if and only if $a=2$. We divide the proof into the following subcases.
\begin{itemize}
\item Let $a=1$ (so that $\bw_1$ has one entry). If $\bw_1^T\be_u\neq 0$, then $\bmu$ has exactly two nonzero entries if and only if $\bw_2$ has exactly one nonzero entry. Equivalently, one of the following conditions holds.
\begin{itemize}
\item $b\geq 1$ and $s$ satisfies (\ref{eqkmn3}), in which case $\bmu=-\be_u+s\be_v$. This establishes 2c.
\item $b=2$ and $s$ satisfies (\ref{eqkmn4}), in which case $\bmu=\left(\cos( \tau\sqrt{2})+\ii \frac{s}{\sqrt{2}}\sin(\tau\sqrt{2})\right)\be_u -s\be_w$, where $B_2=\{v,w\}$. This proves 2d.
\end{itemize}
If $\bw_1^T\be_u=0$ (so that (\ref{eqkmn2}) is satisfied), then $\bmu$ has exactly two nonzero entries if and only if $\bw_2$ has exactly two nonzero entries. Equivalently, one of the following holds.
\begin{itemize}
\item $b=2$ and $\bmu=\left(\ii\frac{1}{\sqrt{2}}\sin(\tau\sqrt{2})+\frac{s}{2}(\cos(\tau\sqrt{2})+1)\right)\be_v+\left(\ii\frac{1}{\sqrt{2}}\sin(\tau\sqrt{2})+\frac{s}{2}(\cos(\tau\sqrt{2})-1)\right)\be_w$ where $B_2=\{v,w\}$. In this case, (\ref{eqkmn2}) simplifies to $s=-\sqrt{2}\cot(\tau\sqrt{2})$, so 2e holds.
%$s+\left(\ii\sin(\tau\sqrt{ab})\frac{1}{\sqrt{ab}}+s(\cos(\tau\sqrt{ab})-1)\frac{1}{2}\right)$
\item $b=3$ and $s$ satisfies (\ref{eqkmn4}), in which case 
$\bmu=-s(\be_w+\be_x)$ where $B_2=\{v,w,x\}$. As $s$ satisfies (\ref{eqkmn2}) and (\ref{eqkmn4}) and $(a,b)=(1,3)$, we get $\tau\in\big\{\frac{2\pi(3k+\ell)}{3\sqrt{3}}:\ell=1,2\ \text{and}\ k\in\Z\big\}$. In particular, $s=-\ii$ if $\tau=\frac{2\pi(3k+1)}{3\sqrt{3}}$ and $s=\ii$ otherwise. This yields 2f, in which case $s$ is a phase factor.
%\left(\ii\sin(\tau\sqrt{3})\frac{1}{\sqrt{3}}+s(\cos(\tau\sqrt{3})-1)\frac{1}{3}\right)
\end{itemize}
\item Let $a\geq 2$ with $u,w\in B_1$ such that $u\neq w$. Since $s$ cannot simultaneously satisfy (\ref{eqkmn3}) and (\ref{eqkmn4}), it cannot happen that $\bw_1^T\be_u\neq 0$ and $\bw_1^T\be_w\neq 0$. Similarly, since $s$ also cannot satisfy both (\ref{eqkmn1}) and (\ref{eqkmn2}), it cannot happen that $\bw_1^T\be_u=\bw_1^T\be_w=0$. Now, suppose $\bw_1^T\be_u=0$ and $\bw_1^T\be_w\neq 0$ (equivalently, (\ref{eqkmn2}) holds). Then $\bmu$ has exactly two nonzero entries if and only if $a=2$ and $\bw_2$ has exactly one nonzero entry. If $b\geq 1$ and $s$ satisfies (\ref{eqkmn3}), then $\bmu=-\be_w+s\be_v$ has exactly two nonzero entries. This proves 2g. Now, if $b\geq 1$ and $s$ does not satisfy (\ref{eqkmn3}), then $\bmu$ has exactly two nonzero entries if and only if $b=1$ or $b=2$ and $s$ satisfies (\ref{eqkmn4}). The case $b=1$ is equivalent to 2d. The case $b=2$ cannot happen because $a=b=2$ contradicts one of (\ref{eqkmn2}) and (\ref{eqkmn4}).
The case when $\bw_1^T\be_u\neq 0$ and $\bw_1^T\be_w=0$ may be shown to be equivalent to 2g with the roles of $B_1$ and $B_2$ reversed.
\end{itemize}
Combining these two subcases completes the proof of the second case. This proves the forward implication. The converse is immediate.
\end{proof}

\end{document}
